\cleardoublepage
\phantomsection
\addcontentsline{toc}{chapter}{はじめに}
\chapter*{はじめに}
\markboth{はじめに}{はじめに}

自然科学のデータ解析では、Python の文法知識だけでは十分ではない。ファイルと実行環境の確認、適切な型によるデータの読込み、計算結果の検証、図による要約までが一連の作業となる。データ量の増加時には、結果の同一性を保ったまま、処理時間とメモリ使用量を抑える工夫も必要である。本書の目的は、読者がこの解析手順を自力で組み立てられるようになることである。

本書の前半では、シェルによるファイル操作、Python の導入と仮想環境、Python の基本構文、データの入出力を扱う。シェルと環境管理は導入時だけの準備ではなく、入力の確認、スクリプトの実行、ログの保存、別の計算機への転送を含む解析基盤である。

中盤の主題は、NumPy、pandas、Polars、Matplotlib、SciPy、Astropy による科学データの処理である。NumPy の役割は数値配列の計算、pandas と Polars の役割は表形式データの整形と集計、Matplotlib の役割は結果の可視化にある。SciPy は統計、フィット、補間、信号処理などの数値計算を、Astropy は物理単位、時刻系、天球座標、FITS 形式を担う。各ライブラリを孤立した道具としてではなく、一つの解析を分担する構成要素として位置付ける。

後半では、反復利用する処理のモジュール化とパッケージ化に加え、テストと依存関係を含む再利用可能な構成を扱う。続く性能改善の章では、計算量、ベクトル化、メモリ使用量、逐次処理、プロファイリングが判断基準となる。巻末の総合課題は、複数の観測装置が \(250\,\mathrm{ns}\) 以内に記録したイベントを同時観測とみなし、ここまでの内容を一つの解析へ統合する課題である。

コードの正常終了は、解析結果の正しさを保証しない。本書が重視する検証は、既知の答えを持つ小規模入力、配列の形状と型、図と残差、高速化前後の結果の一致である。自然科学の解析では、得られた値だけでなく、その値を導いた条件と手順まで説明できなければならない。

Python の言語仕様と標準ライブラリには、公式チュートリアル~\cite{python_tutorial} と公式ライブラリリファレンス~\cite{python_library} を用い、各ライブラリの公式資料とともに巻末へまとめた。
